\uuid{XjUq}
\exo7id{1461}
\titre{exo7 1461}
\auteur{barraud}
\organisation{exo7}
\datecreate{2003-09-01}
\isIndication{false}
\isCorrection{false}
\chapitre{Espace euclidien, espace normé}
\sousChapitre{Produit scalaire, norme}
\cpge{mpsi-mp2i.espaces-prehilbertiens.b, pcsi.espaces-prehilbertiens.b, ptsi.geometrie-espace.b, mpsi-mp2i.integration.c, pcsi.integration.b, ptsi.integration.b, bcpst1.integration-segment.b, pt.structure-prehilbertienne.b, bcpst2.produit-scalaire.a}
\module{Algèbre}
\niveau{L2}
\difficulte{}

\contenu{
\question{
Rappeler l'énoncé de l'inégalité de Cauchy-Schwarz.

Montrer que pour toute fonction continue d'un segment $[a,b]$ dans
$\R$, on a 
$$
  \left(\int_{a}^{b}{f(t)dt}\right)^{2}\leq
  (b-a)\int_{a}^{b}{\Big(f(t)\Big)^{2}dt}
$$

Pour quelles fonctions a-t-on l'égalité ?
}
}
